% Rezeptkarte "Hähnchen-Gemüse-Blech mit Zitronenreis" (4 Portionen). German, 2 pages.
% Features: Kopfleiste mit Portionen/Zeiten, Zutatenliste mit Mengen in drei Gruppen, nummerierte Schritte mit Zeiten und
% Temperaturen, tikz-Ablaufplan (was läuft parallel, mit Beschriftung), Tipps.
\documentclass[11pt,a4paper]{article}
\usepackage[T1]{fontenc}
\usepackage[utf8]{inputenc}
\usepackage{palatino}
\usepackage[scaled=0.9]{helvet}
\usepackage[ngerman]{babel}
\usepackage[margin=2.0cm,top=2.0cm,bottom=2.0cm]{geometry}
\usepackage{xcolor}
\usepackage{tikz}
\usetikzlibrary{calc}
\usepackage{tcolorbox}
\usepackage{enumitem}
\usepackage{booktabs}
\usepackage{array}
\usepackage{caption}
\usepackage{fancyhdr}
\usepackage[locale=DE]{siunitx}
\usepackage{microtype}

\definecolor{tomate}{HTML}{C0392B}
\definecolor{blatt}{HTML}{4F7D3A}
\definecolor{creme}{HTML}{FBF5E6}
\setlength{\parindent}{0pt}
\setlength{\parskip}{0.4em}
\captionsetup{font={small,sf}, labelfont={bf,sf}, labelsep=period, justification=raggedright, singlelinecheck=false}

\pagestyle{fancy}
\fancyhf{}
\renewcommand{\headrulewidth}{0pt}
\fancyfoot[L]{\sffamily\scriptsize Rezeptkarte Nr.\ 17 \textperiodcentered{} Eigenes Rezept für Testzwecke}
\fancyfoot[R]{\sffamily\scriptsize Seite \thepage{} von 2}

\newcommand{\gruppe}[1]{\par\vspace{0.5em}{\sffamily\bfseries\small\color{blatt}#1}\par\vspace{0.15em}}

\begin{document}

% ---- Kopf ----------------------------------------------------------------------------------
{\sffamily\small\color{tomate}\textbf{HAUPTGERICHT \textperiodcentered{} EIN BLECH}}\\[0.5em]
{\fontsize{26}{30}\selectfont\bfseries Hähnchen-Gemüse-Blech\\[-0.1em] mit Zitronenreis}\\[0.3em]
{\itshape Knusprig gebackenes Gemüse, saftiges Hähnchen und ein frischer Joghurt-Dip: fast alles gart gleichzeitig im
Ofen.}

\vspace{0.4em}
\begin{tcolorbox}[colback=creme, colframe=tomate!70, boxrule=0.6pt, arc=2pt, left=6pt, right=6pt, top=4pt, bottom=4pt]
\sffamily\small
\begin{tabular}{@{}l l@{\hspace{2.2em}}l l@{\hspace{2.2em}}l l@{}}
  \textbf{Portionen:} & 4 & \textbf{Aktive Zeit:} & 25 Minuten & \textbf{Schwierigkeit:} & einfach \\
  \textbf{Gesamtzeit:} & 50 Minuten & \textbf{Backofen:} & \qty{220}{\celsius} (Umluft \qty{200}{\celsius}) & \textbf{Pro Portion:} & ca.\ 620 kcal \\
\end{tabular}
\end{tcolorbox}

% ---- Zutaten und Zubereitung nebeneinander ------------------------------------------------------
\vspace{0.4em}
\begin{minipage}[t]{0.37\linewidth}\vspace{0pt}
{\sffamily\Large\bfseries\color{tomate} Zutaten}\\[-0.2em]
{\sffamily\small für 4 Personen}

\gruppe{Für das Blech}
\sffamily\small
\begin{tabular}{@{}r l@{}}
  600 g & Hähnchenbrustfilet \\
  400 g & Süßkartoffeln \\
  2 & rote Paprika (ca.\ 300 g) \\
  1 & Zucchini (ca.\ 250 g) \\
  1 & rote Zwiebel \\
  3 EL & Olivenöl \\
  2 TL & Paprikapulver, edelsüß \\
  1 TL & getrockneter Thymian \\
  1½ TL & Salz \\
  ½ TL & schwarzer Pfeffer \\
\end{tabular}

\gruppe{Für den Zitronenreis}
\sffamily\small
\begin{tabular}{@{}r l@{}}
  250 g & Basmatireis \\
  500 ml & Wasser \\
  ½ TL & Salz \\
  1 & Bio-Zitrone (Schale und Saft) \\
\end{tabular}

\gruppe{Für den Joghurt-Dip}
\sffamily\small
\begin{tabular}{@{}r l@{}}
  200 g & Naturjoghurt \\
  1 & Knoblauchzehe \\
  1 EL & Zitronensaft \\
  1 Bund & Petersilie \\
  1 Prise & Salz \\
\end{tabular}
\end{minipage}\hfill
\begin{minipage}[t]{0.58\linewidth}\vspace{0pt}
{\sffamily\Large\bfseries\color{tomate} Zubereitung}
\begin{enumerate}[leftmargin=2.0em, itemsep=0.45em, label=\textcolor{tomate}{\sffamily\bfseries\arabic*}, labelsep=0.6em]
  \item \textbf{Ofen vorheizen} (Minute 0). Backofen auf \qty{220}{\celsius} Ober-/Unterhitze (Umluft \qty{200}{\celsius})
        vorheizen und ein Backblech mit Backpapier auslegen.
  \item \textbf{Schneiden} (Minute 0 bis 10). Süßkartoffeln schälen und in 2-cm-Würfel schneiden. Paprika und Zucchini
        in 3-cm-Stücke schneiden, die Zwiebel in Spalten. Das Hähnchenfilet in 3-cm-Würfel schneiden.
  \item \textbf{Süßkartoffeln vorgaren} (Minute 10 bis 20). Süßkartoffeln mit 1 EL Öl und ½ TL Salz mischen, auf dem
        Blech verteilen und \qty{10}{min} bei \qty{220}{\celsius} backen.
  \item \textbf{Würzen} (Minute 10 bis 20). Währenddessen Hähnchen, Paprika, Zucchini und Zwiebel mit dem restlichen Öl,
        Paprikapulver, Thymian, 1 TL Salz und dem Pfeffer mischen.
  \item \textbf{Alles auf das Blech} (Minute 20 bis 45). Die gewürzte Mischung zu den Süßkartoffeln geben und weitere
        \qty{25}{min} backen. Nach \qty{12}{min} einmal wenden. Das Hähnchen ist gar, wenn die Kerntemperatur mindestens
        \qty{74}{\celsius} erreicht.
  \item \textbf{Reis kochen} (Minute 28 bis 45). Reis in einem Sieb waschen. Mit 500 ml Wasser und ½ TL Salz aufkochen,
        zugedeckt \qty{12}{min} bei kleinster Hitze quellen lassen und dann \qty{5}{min} ohne Hitze ziehen lassen.
  \item \textbf{Joghurt-Dip} (Minute 30 bis 35). Joghurt mit gepresstem Knoblauch, Zitronensaft und Salz verrühren.
        Petersilie hacken, die Hälfte unterrühren.
  \item \textbf{Servieren} (Minute 45 bis 50). Zitronenschale und Zitronensaft unter den Reis heben. Blech, Reis und Dip
        anrichten und mit der restlichen Petersilie bestreuen.
\end{enumerate}
\end{minipage}

\newpage
% ---- Ablaufplan --------------------------------------------------------------------------------------------
{\sffamily\Large\bfseries\color{tomate} Ablaufplan: Was läuft wann?}\\[0.2em]
{\small Das Blech braucht keine Aufsicht. Nutzen Sie die Zeit im Ofen für Reis und Dip.}

\begin{center}
\begin{tikzpicture}[x=0.255cm, y=0.62cm, font=\sffamily\scriptsize,
    bar/.style={draw=black!60, rounded corners=1pt, minimum height=0.46cm, inner sep=1pt, anchor=west}]
  % rows (top to bottom): Schneiden, Ofen, Reis, Dip, Anrichten
  \node[anchor=east] at (-0.4,0) {Schneiden und würzen};
  \node[anchor=east] at (-0.4,-1) {Backofen};
  \node[anchor=east] at (-0.4,-2) {Reis};
  \node[anchor=east] at (-0.4,-3) {Joghurt-Dip};
  \node[anchor=east] at (-0.4,-4) {Anrichten};
  % bars (minute from, length)
  \node[bar, fill=blatt!35, minimum width=0.255cm*10] at (0,0) {Schneiden};
  \node[bar, fill=blatt!35, minimum width=0.255cm*10] at (10,0) {Würzen};
  \node[bar, fill=orange!25, minimum width=0.255cm*10] at (0,-1) {Vorheizen \qty{220}{\celsius}};
  \node[bar, fill=orange!55, minimum width=0.255cm*10] at (10,-1) {Süßkartoffeln};
  \node[bar, fill=red!40, minimum width=0.255cm*25] at (20,-1) {ganzes Blech, \qty{25}{min}};
  \node[bar, fill=yellow!40, minimum width=0.255cm*12] at (28,-2) {quellen, 12 min};
  \node[bar, fill=yellow!70, minimum width=0.255cm*5] at (40,-2) {ziehen};
  \node[bar, fill=blue!15, minimum width=0.255cm*5] at (30,-3) {rühren};
  \node[bar, fill=tomate!30, minimum width=0.255cm*5] at (45,-4) {Servieren};
  % time axis
  \draw[-latex] (0,-4.8) -- (51.5,-4.8);
  \foreach \m in {0,5,...,50} {\draw (\m,-4.8) -- (\m,-4.95); \node[anchor=north] at (\m,-4.95) {\m};}
  \node[anchor=west] at (51.8,-4.8) {Minuten};
  % marker: wenden
  \draw[dashed, black!60] (32,0.5) -- (32,-4.8);
  \node[anchor=south, fill=white, inner sep=1pt] at (32,0.45) {Blech wenden (Minute 32)};
\end{tikzpicture}
\captionof{figure}{Zeitplan für 50 Minuten. Das Blech ist von Minute 10 bis 45 im Ofen; Reis und Dip entstehen parallel.}
\end{center}

% ---- Naehrwerte --------------------------------------------------------------------------------------------------
\subsection*{\sffamily\large\bfseries\color{tomate} Nährwerte pro Portion (ca.)}
\begin{center}
\sffamily\small
\begin{tabular}{@{}l r r@{}}
  \toprule
  \textbf{Nährwert} & \textbf{Menge} & \textbf{Anteil am Tagesbedarf*} \\
  \midrule
  Energie & 620 kcal & 31 \% \\
  Eiweiß & 44 g & 88 \% \\
  Kohlenhydrate & 78 g & 30 \% \\
  Fett & 14 g & 20 \% \\
  \bottomrule
\end{tabular}\\[0.2em]
{\scriptsize *bezogen auf einen Erwachsenen mit 2\,000 kcal Tagesbedarf; Richtwerte, keine Ernährungsberatung}
\end{center}

% ---- Tipps -----------------------------------------------------------------------------------------------------
\begin{tcolorbox}[colback=creme, colframe=blatt, boxrule=0.6pt, arc=2pt, left=8pt, right=8pt, top=5pt, bottom=5pt,
  title={\sffamily\bfseries Tipps}, fonttitle=\sffamily, coltitle=white, colbacktitle=blatt]
\begin{itemize}[leftmargin=1.3em, itemsep=0.2em]
  \item \textbf{Knusprig statt matschig:} Das Blech nicht überladen. Gemüse und Fleisch sollen in einer Schicht liegen;
        sonst dämpft es statt zu rösten. Im Zweifel zwei Bleche verwenden.
  \item \textbf{Gleichmäßig garen:} Alle Stücke etwa gleich groß schneiden, die Süßkartoffeln etwas kleiner (2 cm) als
        Paprika und Zucchini (3 cm).
  \item \textbf{Sicher gegart:} Das Hähnchen ist durchgegart, wenn die Kerntemperatur mindestens \qty{74}{\celsius} beträgt
        und der Saft klar austritt.
  \item \textbf{Vegetarisch:} Das Hähnchen durch 400 g Kichererbsen (Abtropfgewicht) und 200 g Halloumi ersetzen. Die
        Kichererbsen kommen erst nach Minute 20 aufs Blech.
  \item \textbf{Reste:} Abgekühlt im Kühlschrank bis zu 2 Tage haltbar. Reis nur einmal aufwärmen und gut durcherhitzen.
\end{itemize}
\end{tcolorbox}

\end{document}
